% Замените поясняющий текст собственными результатами.
% В шаблоне нет готового решения и обязательных внешних изображений.
\reportelement{Введение}

Целью работы является перенос запуска приложения, PostgreSQL и Redis в Docker Compose и проба масштабирования; персональные параметры указаны в таблице~\ref{tab:parameters}.

\begin{table}[H]
  \caption{Персональные параметры}\label{tab:parameters}
  \centering
  \begin{tabularx}{\textwidth}{|p{0.46\textwidth}|X|}
    \hline
    Параметр & Значение\\ \hline
    Номер ИСУ & \StudentISU\\ \hline
    Персональный порт (PostgreSQL) & 5489\\ \hline
    Персональный порт (Redis) & 6429\\ \hline
    Персональный порт (приложение-ассистент) & 8059\\ \hline
    Персональный префикс & matveev\\ \hline
    Секретный ключ & secret59matveev\\ \hline
    Имя базы данных & matveev59\\ \hline
  \end{tabularx}
\end{table}

Результат хранится в ветке \texttt{lab-\LabNumber} форка:
\url{\RepositoryURL}. 

\chapter{Ход выполнения}
\section{Подготовка проекта для Docker Compose}

В корне проекта создал файл .env. Он будет содержать все персональные переменные, которые docker-compose подставит в сервисы автоматически. Он изображен на рисунке~\ref{fig:1}.

\reportscreenshot{Screenshot_1.png}{Файл}{fig:1}

Создал файл docker-compose.yml в корне проекта, изображенный на рисунке~\ref{fig:2}. На данном этапе описаны два сервиса: PostgreSQL (db) и приложение (app). Использованы переменные из .env.

\reportscreenshot{Screenshot_2.png}{Файл}{fig:2}

Запустил все сервисы и убедился, что оба контейнера запущены и приложение отвечает на рисунке~\ref{fig:3}.

\reportscreenshot{Screenshot_3.png}{Проверка}{fig:3}

\section{Добавление Redis для кэширования погоды}

Добавил новый сервис redis в docker-compose.yml. Добавил зависимость app от Redis. Новый файл изображен на рисунках~\ref{fig:4} и~\ref{fig:5}.

\reportscreenshot{Screenshot_4.png}{Файл}{fig:4}
\reportscreenshot{Screenshot_5.png}{Файл}{fig:5}

Интегрировал в main.py redis, добавлен эндпоинт /weather/{city}, имеет логику кэширования, добавлен эндпоинт /cache/stats. Новый код изображен на рисунках~\ref{fig:6},~\ref{fig:7} и~\ref{fig:8}.

\reportscreenshot{Screenshot_6.png}{Файл}{fig:6}
\reportscreenshot{Screenshot_7.png}{Файл}{fig:7}
\reportscreenshot{Screenshot_8.png}{Файл}{fig:8}

\section{Запуск и проверка полного стека}

Пересобрал и запустил все сервисы. Проверил порядок запуска на рисунке~\ref{fig:9}.

\reportscreenshot{Screenshot_9.png}{Проверка}{fig:9}

Проверил работу кэширования на рисунке~\ref{fig:10}.

\reportscreenshot{Screenshot_10.png}{Проверка}{fig:10}

Исследовал networks и volumes на рисунке~\ref{fig:11}.

\reportscreenshot{Screenshot_11.png}{Команды и выводы}{fig:11}

\section{Масштабирование сервиса}

Создал отдельный файл docker-compose.scale.yml. Для масштабирования убрал \texttt{container-name} у app и изменил порт на диапазон. Файл изображен на рисунке~\ref{fig:12}.

\reportscreenshot{Screenshot_12.png}{Код}{fig:12}

Запустил с масштабированием (3 экземпляра app) и проверил на рисунке~\ref{fig:13}.

\reportscreenshot{Screenshot_13.png}{Проверка}{fig:13}

Проверил работу масштабированных экземпляров на рисунке~\ref{fig:14}.

\reportscreenshot{Screenshot_14.png}{Проверка}{fig:14}

\section{Финальный артефакт}

Финайльный артефакт изображен на рисунке~\ref{fig:15}.

\reportscreenshot{Screenshot_15.png}{Артефакт}{fig:15}

\reportelement{Заключение}

В ходе выполнения лабораторной работы №4 были изучены возможности Docker Compose для оркестрации многоконтейнерных приложений. Была создана конфигурация, обеспечивающая совместную работу приложения, базы данных PostgreSQL и Redis.

В процессе работы была реализована интеграция Redis для кэширования результатов запросов о погоде и зависимости между сервисами (depends-on). Также было выполнено масштабирование приложения до трёх экземпляров и проверена их совместная работа с общей базой данных и кэшем.


% Оставьте только использованные источники и цитируйте их в тексте.
\begin{thebibliography}{9}
  \bibitem{docker-docs}
  Docker Documentation: документация Docker.
  URL: \url{https://docs.docker.com/}
  (дата обращения: укажите свою дату).
\end{thebibliography}

